\documentclass[sigconf,nonacm]{acmart}
\usepackage{amsmath,mathtools,array}
\settopmatter{printfolios=true,printacmref=false}
\setcopyright{none}
\renewcommand{\shortauthors}{}
\setlength{\emergencystretch}{0.7em}
\newcommand{\E}{\mathbb E}
\newcommand{\R}{\mathbb R}
\newcommand{\TV}{\operatorname{TV}}
\newcommand{\Hell}{\mathsf h}
\newcommand{\U}{\mathcal U}
\newcommand{\A}{\mathcal A}
\newcommand{\Obs}{\mathcal O}
\newcommand{\K}{\mathcal K}
\newcommand{\B}{\mathcal B}
\newcommand{\I}{\mathcal I}
\newcommand{\cE}{\mathcal E}
\newcommand{\wh}{\widehat}
\newcommand{\eps}{\varepsilon}
\newcommand{\rank}{\operatorname{rank}}
\newcommand{\tr}{\operatorname{tr}}
\newcommand{\diag}{\operatorname{diag}}
\newcommand{\softO}{\widetilde O}
\newcommand{\st}{\star}
\newcommand{\fused}[1]{\lVert #1\rVert_{\mathrm F}}
\newcommand{\dual}[1]{\lVert #1\rVert_{\mathrm F,*}}
\newcommand{\pn}[1]{\lVert #1\rVert_{\Pi}}
\newcommand{\ip}[2]{\langle #1,#2\rangle}
\newcommand{\sect}[1]{Section~\ref{#1}}
\begin{document}
\title[Learning PSRs with Ordinary MLE]{Learning Predictive State Representations with Ordinary Maximum Likelihood and Model Planning}
\begin{abstract}
How can a learner decide whether to trust a single maximum likelihood fit of a partially observed system? We construct a trajectory reward from the fitted probabilities of observable tests. Its expected value bounds the trajectory prediction error of each policy. When planning finds a large reward, the optimizing coefficients specify actions and scores for testing the uncertain prediction in the environment. Sampling the resulting experiment either adds coverage of a true predictive direction or exposes a discrepancy that later fits must account for. For a finite realizable class of uniformly B-stable predictive state representations, the algorithm learns an $\eps$-optimal policy in $\softO(r\Lambda_B^4 A U^2 H^4\beta/\eps^2)$ episodes with $\softO(rH)$ calls each to ordinary maximum likelihood estimation and single-model planning. Here $r$ is the true reachable rank, $U$ counts core action strings, $A$ is the number of actions, and $\beta$ measures likelihood complexity. The final model supports planning for any bounded observable reward.
\end{abstract}
\maketitle

\section{Introduction}
Suppose a learner can fit a model of a partially observed system by ordinary maximum likelihood and plan in that model. How many times must it call these routines before it can trust its predictions under every policy? Between calls it may collect many episodes using the same fit. To justify another call, it needs an experiment whose information will still count after the model changes.

Consider a job that arrives with a positive or negative cue. The agent routes the job to one of two stations, sees the same ready signal at either station, and chooses a processing mode. In \sect{sec:example}, two models agree on everything observed when the agent always uses mode $b$. Yet after a negative cue and a right turn, one model assigns probability one to success in mode $a$, and the other assigns probability zero. A predictive state representation (PSR) encodes an observed history by the probabilities of future tests \cite{littman,singh}. The test that selects $a$ and records success has the same meaning under both fits. Its coordinate at the very same history can move from one to zero. Even old training rows change when the learner evaluates them in the next fit. Growth of a covariance matrix computed in today's model therefore gives no immediate bound on the number of future fits.

We use fitted predictions to choose an experiment the learner can run. At a new history, the learner finds a linear combination of test probabilities that is large there and small on the training histories. Its coefficients specify actions to take and a score to record. When the resulting exploration reward is large, a new batch either reaches a predictive direction poorly covered in the true system, or reveals a difference between the fitted and true laws of the chosen experiment. In the latter case we interrupt the experiment at a random stage and test its local predictions. Every subsequent maximum likelihood fit must account for those observations, even if it assigns new coordinates to every stored history. The first alternative increases empirical coverage of true predictive states. The second increases coverage under the distributions of the diagnostic policies. Both increases are measured in the fixed predictive space of the true system.

For a finite realizable class of uniformly B-stable PSRs, this argument gives $\softO(rH)$ calls each to ordinary MLE and single-model planning, where $r$ is the reachable rank of the true system and $H$ the horizon. With probability at least $1-\delta$, the final fit approximates the trajectory law of every policy to total variation (TV) error at most $\eps/4$. One planning call then gives an $\eps$-optimal policy for any specified observable reward in $[0,1]$. The real-episode bound is $\softO(r\Lambda_B^4 A U^2 H^4\beta/\eps^2)$. Here $\Lambda_B$ controls the propagation of prediction error, $A$ counts actions, $U$ counts distinct core action strings, and $\beta$ measures likelihood complexity.

\paragraph{From model disagreement to an executable test.}
Let $P_\theta^\pi$ denote the law of the observed trajectory under model $\theta$ and policy $\pi$. Reward-free OMLE \cite[Algorithm~2, line~3]{omle} searches its likelihood confidence set $\mathcal C_k$ for a policy and two models attaining
\[
 \max_{\pi}\max_{\theta,\theta'\in\mathcal C_k}\TV(P_\theta^\pi,P_{\theta'}^\pi).
\]
All-Policy Model-Estimation E2D achieves uniform trajectory accuracy for B-stable PSRs by optimizing distributions over exploration policies and output models \cite[Theorem~H.4]{bstable}. PSR-UCB plans in one fitted model with a prediction bonus and fits under constraints on the probabilities of sampled prefixes \cite{psrucb}. Table~\ref{tab:methods} compares the episode bounds and optimization steps.

Our planner receives a trajectory reward $b_k$ computed from one ordinary fit $\wh\theta_k$ and the logged histories. With high probability,
\begin{equation}
 \TV(P_\st^\pi,P_{\wh\theta_k}^\pi)
 \le \E_{\wh\theta_k,\pi} b_k \qquad\text{for every }\pi .
 \label{eq:intro}
\end{equation}
Maximizing the right side tells us when to stop. If it is large, the optimizing coefficients prescribe the tests for the next batch. Section~\ref{sec:certificate} derives the reward, \sect{sec:diagnostic} turns its coefficients into experiments, and \sect{sec:progress} counts the rounds in which either kind of information accumulates. A round can execute several test controllers between oracle calls. Policy switches and deployment batches are separate measures of adaptivity \cite{xiong}.

\begin{table*}[t]
\caption{Exploration guarantees and optimization steps. The episode column suppresses $\eps^{-2}$, likelihood complexity and logarithms apart from OMLE's operator-dependent $\iota$. Here $D$ is the largest core-test list, $d$ bounds rank throughout the E2D class, and $r$ is the rank parameter in the other rows.}
\label{tab:methods}
\small\centering
\setlength{\tabcolsep}{4pt}\renewcommand{\arraystretch}{1.15}
\begin{tabular}{@{}>{\raggedright\arraybackslash}p{.17\textwidth}>{\raggedright\arraybackslash}p{.21\textwidth}>{\raggedright\arraybackslash}p{.235\textwidth}>{\raggedright\arraybackslash}p{.32\textwidth}@{}}
\toprule
Method & Target and condition & Episode upper bound & Optimization and frequency\\
\midrule
OMLE \cite[Thm.~9]{bstable} & Known-reward policy, B-stability & $rAU\Lambda_B^2 H^2\iota$ & One optimistic model--policy search per $H$ episodes.\\
All-Policy E2D \cite[Thm.~H.4]{bstable} & Uniform TV, B-stability or weak B-stability & $dAU\Lambda_B^2 H^2$ & One minimax update over policy and model distributions per episode.\\
PSR-UCB \cite[Thm.~1]{psrucb} & Uniform TV, signed $\ell_1$ stability & $(r+U^2H/\gamma^2)\,rD$ $\times H^3 A^2 U^4/\gamma^4$ & One prefix-constrained fit and one bonus plan per $H$ episodes.\\
This work & Uniform TV, B-stability & $rAU^2\Lambda_B^4 H^4$ & $\softO(rH)$ calls each to ordinary MLE and single-model planning.\\
\bottomrule
\end{tabular}
\Description{Comparison of learning targets, structural assumptions, episode bounds and optimization frequency.}
\end{table*}

\section{Setting and computational access}\label{sec:setting}
The horizon is $H$. At stage $h$ the environment emits $o_h\in\Obs$ and the learner takes $a_h\in\A$, where both alphabets are finite and $A=|\A|$. Write $\tau_h=(o_1,a_1,\ldots,o_h,a_h)$ and $x=\tau_{h-1}$. A policy specifies $\pi_h(a\mid x,o)$ on all histories. A model $\theta$ and a policy induce a trajectory law $P_\theta^\pi$ and value $V_{\theta,R}^\pi=\E_{\theta,\pi}R(\tau_H)$ for a known observable reward $R\in[0,1]$. A star denotes the true environment.

\paragraph{Predictive coordinates.}
A stage-$h$ test $j=(o_{h:h+\ell},a_{h:h+\ell-1})$ ends with an observation. Its probability is the chance of its observation string when its actions are imposed. All models share known finite lists $\U_h$ of core tests. Write $\alpha(j)$ for a test's action string, $U_h$ for the number of distinct such strings, $U=\max_h U_h$ and $D=\max_h|\U_h|$. The predictive state $q_h^\theta(x)$ lists the test probabilities. Add a terminal core $\U_{H+1}=\{\varnothing\}$ of value one. A proper B-representation specifies an initial vector and linear maps satisfying
\begin{equation}
 B_h^\theta(o,a)q_h^\theta(x)=
 P_\theta(o\mid x)\,q_{h+1}^\theta(xoa).
 \label{eq:update}
\end{equation}
Their products give the observable probabilities of the model and a normalized trajectory law under any policy. The action is chosen after $o$, so it does not affect the factor $P_\theta(o\mid x)$. On a history of zero model probability use a fixed valid predictive state at the same stage as a conditional version, with continuations from that state.

\paragraph{Which dimensions enter the problem?}
A group of core tests with the same action string can be observed in one experiment: draw the string, execute its actions and record the observations. The tests within a group have the same length and distinct observation strings. For example, all single-observation tests use the empty action string. With $d$ possible observations this gives $d$ coordinates measurable by one experimental design. Thus $D$ determines how many coordinates to compute and $U$ determines how many different strings share the sampling budget. A third dimension enters when we ask how many of the resulting vectors are linearly independent as the history changes.

\paragraph{Reachable rank and stability.}
The finite model class $\Theta$ contains the truth. At every stage its true reachable vectors $q_h^\st(x)$ span a space of dimension at most $r$, with $r\ge1$. To describe propagation of signed errors, let $\mathcal J_h$ contain the core tests maximal under the prefix order. For a signed vector $v$, put
\begin{align}
 N_{12}(v)&=\Big(\sum_\alpha\big(\sum_{j:\alpha(j)=\alpha}|v_j|\big)^2\Big)^{1/2},
 \nonumber\\
 N_{\Pi}(v)&=\max_\pi\sum_{j\in\mathcal J_h}\pi(j)|v_j|,
 &\fused{v}&=\max\{N_{12}(v),N_\Pi(v)\}.
 \label{eq:fused}
\end{align}
Here $\pi(j)$ is the product of action probabilities along the test. For an array of full suffixes, $\pn{\cdot}$ is the maximum policy-weighted absolute sum over those suffixes. Let $\B^\theta_{H:h}$ map $v$ to the array obtained by applying every full product $B^\theta_H\cdots B^\theta_h$. Every candidate satisfies the signed-input B-stability inequality \cite{bstable}
\begin{equation}
 \pn{\B^\theta_{H:h}v}\le \Lambda_B\fused{v}.
 \label{eq:stable}
\end{equation}
Take $\Lambda_B\ge1$. Each action-string group has probability mass at most one, and the maximal leaves reached by a policy are disjoint. Thus $\fused{q_h^\theta(x)}\le\sqrt U$. We use the dual norm $\dual{c}=\sup_{\fused v\le1}c^\top v$ at its indicated stage.

\paragraph{Why stability is stated on signed inputs.}
A proper state has nonnegative test probabilities, but the difference of two states has both signs. A local change in dynamics creates a second signed vector, $(B_h^\theta-B_h^{\theta'})q$. Either vector has to be propagated through the remaining steps before its effect on reward is known. The left side of \eqref{eq:stable} measures exactly that effect after maximizing over continuations. On the right, $N_{12}$ groups outcomes observable under a common action string. The second term allows an experiment to choose an action after seeing an observation. If one branch calls for action $a$ and another for action $b$, one contingent policy can implement both prescriptions. The fused norm retains these two ways of testing a prediction. Its dual will specify a random experiment that executes their mixture.

\paragraph{Access model.}
An episode controller $\nu_i$ is chosen before episode $i$ and may adapt its actions to observations within the episode. Its log $\omega_i$ includes the choices of experimental branches. The ordinary cumulative MLE oracle returns a model in $\Theta$ with total log likelihood within $\xi$ of $\sup_{\theta\in\Theta}\sum_i\log p_\theta^{\nu_i}(\omega_i)$. The planning oracle takes one model, a bounded evaluable history reward, and a tolerance, and returns an executable policy within that tolerance of the optimal value. We assume polynomial-size oracle outputs, efficient execution and simulation, and polynomial-bit evaluation of likelihoods and predictive features to prescribed precision. The learner has the stated structural parameter bounds. Computational costs outside the oracles are measured in the explicit representation, alphabet sizes, data size and accuracy parameters.

\section{Algorithm and main result}\label{sec:algorithm}
For the finite-class result write
\begin{equation}
 \begin{gathered}
 L=H+1,\qquad W=\Lambda_B\sqrt U,\qquad Y=\Lambda_B^2 A U,\\
 \beta=\max\{1,\log(3|\Theta|/\delta)+\xi/2\}.
 \end{gathered}
 \label{eq:params}
\end{equation}
Here $W$ bounds a stage score, while $Y$ converts the squared Hellinger error of a randomized test into squared local prediction error.

\begin{theorem}\label{thm:main}
Under the assumptions in \sect{sec:setting}, for $0<\eps\le1$ and $0<\delta<1$, the algorithm below returns, with probability at least $1-\delta$, a policy $\wh\pi$ satisfying $\sup_\pi V_{\st,R}^\pi-V_{\st,R}^{\wh\pi}\le\eps$. Its real-episode budget $T$ and number of calls to each oracle are, respectively,
\[
 T=\softO\!\left(\frac{r\Lambda_B^4 A U^2 H^4\beta}{\eps^2}\right),
 \qquad K_{\rm oracle}=\softO(rH).
\]
The remaining computation is polynomial under the stated access model.
\end{theorem}

\begin{corollary}[Planning after exploration]\label{cor:reward}
On the success event, the final fit satisfies $\sup_\pi\TV(P_\st^\pi,P_{\wh\theta_k}^\pi)\le\eps/4$. The same fit supports an $\eps$-optimal policy for each observable $R\in[0,1]$ with one additional planning call at accuracy $\eps/2$ for each reward.
\end{corollary}
To see the consequence of uniformity, let $\pi^\star$ be optimal under the truth, $\hat\pi$ the planner's policy, and $\Delta=\sup_\pi\TV(P_\st^\pi,P_{\wh\theta_k}^\pi)$. In order from left to right, transfer the optimum's value, plan in the fit, and transfer the output's value back:
\[
 V_{\st,R}^{\pi^\star}-V_{\st,R}^{\hat\pi}
 \le \Delta+\eps/2+\Delta .
\]
Uniformity covers both policies in this inequality.

\subsection{Exploration with ordinary MLE}
Store a data set $\mathcal D$ of complete logged episodes, a multiset $\I_h$ of training prefixes at each stage, and the number $N_0$ of initial-core episodes. The log includes both the physical trajectory and the randomized choices used to collect it. The stage buffers can therefore be reconstructed from $\mathcal D$. The learner reevaluates the predictive features of these same histories whenever the fit changes.

\begin{figure}[t]
\hrule\smallskip
\textbf{Algorithm 1}\quad Learning with a fitted PSR\par
\smallskip
\emph{Input.} Model class, core tests, structural bounds, accuracy $\eps$, confidence $1-\delta$, and task reward $R$.\par
\smallskip
Initialize $\mathcal D$ and the buffers to empty. Until the deterministic cap in Appendix~\ref{app:budget}:
\begin{enumerate}
\item Fit $\wh\theta_k$ by ordinary MLE on $\mathcal D$. Build the stage programs from the reencoded buffers.
\item Plan for the attained reward $\widetilde b_k$ in $\wh\theta_k$ and estimate the returned policy's bonus value by fitted-model simulation. Use tolerance $\eps/128$ in both steps.
\item If the estimate is at most $3\eps/16$, return the policy obtained by planning for $R$ in $\wh\theta_k$ to accuracy $\eps/2$.
\item Freeze $\wh\theta_k$, the policy $\pi^k$, and the score-selection rule. Collect the training and nested diagnostic batches. Append every episode to $\mathcal D$, update the training buffers and $N_0$, then refit.
\end{enumerate}
Return failure if the cap is reached.\par\smallskip\hrule
\caption{The only optimization over models occurs in step 1. The two planning calls optimize different rewards in the same model.}
\Description{Pseudocode for fitting a model, planning to maximize its uncertainty reward, and either returning a policy or collecting more data.}
\label{fig:algorithm}
\end{figure}

Figure~\ref{fig:algorithm} gives the decision rule. The trajectory reward is the sum of a scalar initial-state term and the attained support values of the stage programs in \sect{sec:certificate}. The numerical rule has a fixed tie convention, so revisiting a history in the same round produces the same coefficient and the same experiment. When exploration continues the learner freezes this rule before gathering the first episode of either batch.

\emph{Training batch.} At each $h=1,\ldots,H$, run $m_k$ fresh episodes following $\pi^k$ to $x=\tau_{h-1}$. Flip a fair coin. The first branch draws a uniform stage-$h$ core action string and executes it. The second observes $o_h$, takes a uniform action, and executes a uniform stage-$(h+1)$ core string. The second branch samples the core string in the roll-in and intervention of \citet[Eq.~(10)]{omle}. The two coin outcomes measure the current prediction $q_h$ and the joint prediction $B_h(o,a)q_h$ under the same roll-in distribution. Retain $x$ in the stage buffer $\I_h$. Run $m_k$ additional initial-core experiments from the empty history and increment their count $N_0$. Every experiment uses a common completion rule through the horizon. For a chosen string, one observes which listed observation outcome occurred, or that none occurred.

\emph{Nested diagnostic batch.} The coefficient rule and $\pi^k$ specify the controller $\rho_k$ constructed in \sect{sec:diagnostic}. In each of $n_N$ fresh episodes draw $j$ uniformly from $\{0,\ldots,H-1\}$. The cut $E_0(\rho_k)$ runs an initial uniform-core experiment. For $j\ge1$, $E_j(\rho_k)$ follows $\rho_k$ through stage $j-1$, observes $o_j$, replaces the action by a uniform one, and executes a uniform next-core string. Log the cut, choices already drawn, and new experimental choices. A scheduled choice superseded by the cut receives a null label. The diagnostic data contribute complete likelihoods. Only the designated training histories supply rows to the stage programs.

A frozen round spends $(H+1)m_k+n_N$ real episodes before paying for another MLE call. The $m_k$ episodes at each training stage provide rows at that stage's roll-in distribution. The $n_N$ nested experiments constrain the law induced by the diagnostic as a whole. Appendix~\ref{app:budget} specifies these batch sizes.

\paragraph{Relation to statistical rates.}
For the uniform-TV objective of Corollary~\ref{cor:reward}, the direct comparison is All-Policy E2D \cite[Theorem~H.4]{bstable}. Its rank bound $d$ applies to every candidate. Our $r$ bounds the span of reachable predictions in the true system. At matched ranks and after separating likelihood and logarithmic factors, our episode upper bound has an additional factor $\Lambda_B^2 U H^2$. Section~\ref{sec:progress} traces it to the accuracy needed to test a fitted direction. For the known-reward objective, the OMLE bound in \citet[Theorem~9]{bstable} also includes a logarithm involving the true one-step operator norm. The original OMLE guarantee \cite[Theorem~4.5]{omle} expresses the structural dependence as a polynomial.

For the signed $\ell_1$ continuation bound $\gamma^{-1}$, we may take $\Lambda_B=\sqrt U/\gamma$ since $\|v\|_1\le\sqrt U\fused v$ \cite[Appendix~D.4]{bstable}. This is the full-trajectory inequality for $m_1$ in \citet[Condition~4.3]{omle}. Our episode bound then becomes $\softO(rAU^4H^4\beta/(\gamma^4\eps^2))$. Under this parameterization, the main Theorem~1 of \citet{psrucb} states an episode bound with polynomial and likelihood factors
\[
 \left(r+U^2H/\gamma^2\right)
 \frac{rD H^3 A^2 U^4\beta_{\rm UCB}}{\gamma^4\eps^2}.
\]
After separating the likelihood budgets, the ratio of our polynomial factor to theirs is $H/[AD(r+U^2H/\gamma^2)]$. The dimension $D$ enters our explicit cone programs; the statistical effect of enlarging the candidate class enters through $\beta$.

\section{Two outcomes of the same optimal test}\label{sec:example}
A large exploration reward can lead to two different kinds of data. We exhibit both in a three-stage system, solving the stage program and following its optimizer through the experiments in Algorithm~1.

\subsection{A prediction hidden by the task policy}
A job arrives with a fair binary cue $c\in\{+,-\}$. The first action chooses a route $L$ or $R$. Either station emits the same observation, \emph{ready}. At this point the job has an unobserved type $t\in\{0,1\}$. In the true system it is one exactly after $(+,R)$. A second model assigns type one after either cue on route $R$. Processing in mode $a$ returns the type as the binary final observation $y$. Mode $b$ returns an independent fair bit. The reward is $y$.

In the second model, always routing right and using $a$ attains value one. In the true system that policy earns $1/2$. A true optimum routes right and uses $a$ on $+$, then uses $b$ on $-$, for value $3/4$. An agent choosing $b$ after every route observes exactly the same trajectory law in both systems. Nevertheless the cue and the route retain information about the consequence of changing the processing mode.

Use the two cue outcomes as the stage-one core. At stage two the two tests observe \emph{ready}, impose $a$, and observe $y=0$ or $1$. The final core consists of the two possible outcomes of $y$. Each core uses a single action string, so $U=1$. Write $e_0,e_1$ for the two type vectors and $m=(1/2,1/2)^\top$. In either model $q_1=m$, $q_2=e_t$, and the observable updates on the effective actions are
\begin{align*}
 B_1(c,s)v &= v_c e_{t(c,s)},\\
 B_2(\mathrm{ready},a)v &=v, &
 B_2(\mathrm{ready},b)v &=m\mathbf{1}^{\top}v,\\
 B_3(y,\cdot)v &= v_y .
\end{align*}
Here $s$ denotes the selected route. The two actions are indexed by their stage-specific names. Every basis vector induces a probability law on the remaining outcomes under each policy, so the continuation norm of a signed vector is at most its $\ell_1$ norm. With one action string the fused norm is exactly $\ell_1$. We supply $\Lambda_B=2$ as a common stability upper bound.

\subsection{The stage optimizer and its experiment}
Consider a buffer with $n$ designated training prefixes at each stage and $n$ initial-core episodes. The training roll-in always routes left and uses $b$. Both models encode every stage-two training prefix as $e_0$ and every stage-three training prefix as $m$. The stage-two lifted row at type $i$ is $(e_i,e_i,m)$, listing the current predictions and the two joint blocks. For coefficients $(u,w_a,w_b)$ write its evaluation as
\[
 g_i=u_i+(w_a)_i+w_b^\top m .
\]
The structured ball in \eqref{eq:geometry} imposes $\|u\|_\infty\le1$ and $\|w_a\|_\infty+\|w_b\|_\infty\le1$. The historical constraint in \eqref{eq:program} reduces to $n g_0^2\le S$. At a prefix encoded as $e_1$ the objective is $g_1$. The norm constraints give $g_1\le2$, and the choice
\[
 u=e_1,\qquad w_a=e_1,\qquad w_b=0
\]
attains two with $g_0=0$. Thus the exact stage support is two. Every optimizer satisfies $g_1-g_0\ge 2-\sqrt{S/n}$. For the displayed optimizer, each nonzero coefficient uses a single string with mass one. Both the current-core and the transition experiment take action $a$ after \emph{ready} and score $\mathbf1\{y=1\}$. Each has conditional mean $i$ at type $i$.

\subsection{What one high score can buy}
First let the fit equal the truth. Route right on $+$ and left on $-$, then use $b$. Half the stage-two histories have type one. The displayed coefficient has zero historical energy in the true coordinates, so the expected true support is at least one. Each new stage-two training prefix from the positive cue adds an outer product in the previously unobserved type-one direction. All diagnostic laws agree because the fitted dynamics are correct.

Now use the second model as the fit. Route right on $-$ and left on $+$, again using $b$. On the negative cue the fit predicts $e_1$ where the environment has $e_0$. The expected fitted stage support is at least one, and the policy's trajectory TV is zero. In the true system the root and every queried stage row coincide with their respective training rows. Hence any coefficient tuple of total historical energy at most $\cE$ has summed evaluation at most $\sqrt{(H+1)\cE/n}$ by Cauchy--Schwarz. Taking $n$ sufficiently large makes the true support smaller than the first threshold in \eqref{eq:either}. In contrast, each displayed stage-two diagnostic gives opposite deterministic outcomes on the negative cue. Its TV is at least $1/2$, and the logged mixture in \eqref{eq:rho} has TV at least $1/(3H)$.

Finally, consider cut $2$ of that mixture. With probability $1/3$ the logged component is the full policy. Within it the negative cue occurs with probability $1/2$ and the fresh uniform action selects $a$ with probability $1/2$. On this observable branch the two terminal laws have disjoint support. The branch has mass $1/12$ under both models, so its contribution to squared Hellinger distance is $1/12$. A uniform cut chooses stage two with probability $1/3$, giving at least $1/36$ per nested episode. Consequently $n_N$ repetitions charge at least $n_N/36$ to the likelihood budget of any later fit with these dynamics. The two calculations concern a fixed policy and a designated buffer. The general proof establishes progress for the policy actually returned by the planner, whichever maximizing directions its stage programs select.

\section{Certifying trajectory error}\label{sec:certificate}
Fix a round and abbreviate $\wh\theta_k$ by a hat. We first identify a local error whose expectation is taken under the fitted model. The direction of this expansion determines which features the planner needs. For a continuation operator at $H+1$, use the identity. Define
\begin{align}
 F_0 &= \tfrac12\pn{\B^\st_{H:1}(q_1^\st-\wh q_1)},\nonumber\\
 F_h(q)&=\tfrac12\max_\alpha\sum_{o,a}\alpha(a\mid o)
 \pn{\B^\st_{H:h+1}(B_h^\st(o,a)-\wh B_h(o,a))q}.
 \label{eq:F}
\end{align}
Here $\alpha$ is an action rule following the current observation. Each suffix norm can choose its own continuation, as a history policy can branch after $(o,a)$.

\paragraph{The direction of expansion.}
Subtract the two products for a trajectory by replacing one factor at a time. The summand for stage $h$ has true operators to its left (the future), the difference $B_h^\st-\wh B_h$ in the middle, and fitted operators to its right (the prefix). Acting on the fitted initial vector, the prefix product gives its fitted observation probability times $\wh q_h(x)$. Multiply by the policy's action probabilities and sum absolute errors over trajectories. Summation over the prefix now gives the fitted expectation. Thus
\begin{equation}
 \TV(P_\st^\pi,P_{\wh\theta}^\pi)
 \le F_0+\sum_{h=1}^H\E_{\wh\theta,\pi} F_h(\wh q_h(x)).
 \label{eq:telescope}
\end{equation}
Appendix~\ref{app:product} gives the product identity, including the initial-state term. The future operators in $F_h$ are unknown. We next construct an observable upper bound on this seminorm.

\paragraph{The features seen by an error direction.}
For either model define
\[
 z_h^\theta(o,a;x)=B_h^\theta(o,a)q_h^\theta(x),\qquad
 \phi_h^\theta(x)=\big(q_h^\theta(x),(z_h^\theta(o,a;x))_{o,a}\big).
\]
Write $g_v(\phi)=u^\top q+\sum_{o,a}w_{oa}^\top z_{oa}$ for $v=(u,w)$. Use the stage-dependent coefficient set
\begin{equation}
 \K_h=\left\{(u,w):
 \dual u\le\tfrac{\Lambda_B}{2},\
 \sum_a\dual{w_{oa}}\le\tfrac{\Lambda_B}{2}\;\;(\forall o)\right\}.
 \label{eq:geometry}
\end{equation}
The norms on $u$ and $w$ correspond to stages $h$ and $h+1$. Every such functional has magnitude at most $W$ on valid lifted states by \eqref{eq:update}. Its true lift is a fixed linear map of $q_h^\st$, a fact we will use to count updates.

\subsection{A witness at a new history must be feasible on old histories}
The seminorm $F_h$ depends on both models. The learner can still put its supporting directions inside a known set. We spell out this change of variables because it is what makes optimization in a single fit possible.

Fix the query and a maximizing action rule in \eqref{eq:F}. For each $(o,a)$ choose a suffix policy and signs realizing the absolute sum in that suffix norm. Let $d_{oa}$ be the corresponding signed row of the true continuation map, summed with its policy weights. Thus $|d_{oa}^\top y|\le\pn{\B^\st_{H:h+1}y}$ for every signed $y$. Write
\begin{align*}
 u^\top &=\tfrac12\sum_{o,a}\alpha(a\mid o)d_{oa}^\top B_h^\star(o,a),\\
 w_{oa}^\top &=-\tfrac12\alpha(a\mid o)d_{oa}^\top .
\end{align*}
The first row is a single signed full continuation: at observation $o$ pick $a$ with probability $\alpha(a\mid o)$, and continue with the policy already chosen for that branch. Even though the number of observations can be large, it follows that $|u^\top y|\le\pn{\B^\st_{H:h}y}/2$. The second row has dual norm at most $\alpha(a\mid o)\Lambda_B/2$. Summing over actions at fixed $o$ proves both constraints in \eqref{eq:geometry}.

At the query this choice gives $g_v(\wh\phi_h(x))=F_h(\wh q_h(x))$. At another input the chosen policy and signs are still available to the maximization defining $F_h$. Reversing all signs also remains feasible. Hence
\begin{equation}
 |g_v(\wh\phi_h(x'))|\le F_h(\wh q_h(x'))\quad\text{for every }x'.
 \label{eq:global}
\end{equation}
Equality at the query and domination on the training set do different work. Equality makes the objective large enough to certify the error of interest. Domination implies that the same witness passes a constraint on its accumulated historical energy. The unknown continuation coefficients can now be replaced by the whole structured ball, provided we can estimate how large that energy may be.

\subsection{What the likelihood of a complete episode controls}
A training episode reaches a random prefix and only then chooses the experiment that produces its remaining data. At its observed prefix $X$ the desired error is a discrepancy between suffix laws. The oracle, however, maximizes the joint log likelihood of $(X,Z)$. An observed suffix may favor a candidate that assigned little probability to its prefix, or conversely. The following inequality accounts for both factors in the same likelihood ratio.

\begin{lemma}[Conditional likelihood calibration]\label{lem:mle}
In each episode select a predictable rule for splitting its recorded data into a prefix $X_i$ and a suffix $Z_i$. For a candidate $\theta$, let $a_{\theta,i}(x)=\Hell^2(P_{\st,i}(Z_i\mid x),P_{\theta,i}(Z_i\mid x))$, where $\Hell^2(P,Q)=1-\sum_z\sqrt{P(z)Q(z)}$. With probability at least $1-\delta/3$, at every refit on $n$ episodes,
\[
 \sum_{i\le n} a_{\wh\theta,i}(X_i)\le\beta .
\]
\end{lemma}
The proof in Appendix~\ref{app:likelihood} applies to the ordinary joint likelihood of every complete episode. In training episodes the designated prefix ends immediately before the test coin. In initial-core and nested episodes it is empty. The first choice controls conditional error at the histories actually observed. The second controls a population distance between complete experimental laws.

For a fixed candidate write $p,q$ for the true and candidate joint laws and $a(x)$ for conditional squared Hellinger error. Integrating a square-root likelihood ratio over the suffix yields
\[
 \sum_z p(x,z)\sqrt{\frac{q(x,z)}{p(x,z)}}
   =\sqrt{p_X(x)q_X(x)}\,\bigl(1-a(x)\bigr).
\]
Multiply by $e^{a(x)}$. Since $e^t(1-t)\le1$ on $[0,1]$, summing over $x$ gives at most one, even as the experiment changes with past data. Iterating this calculation and applying the MLE likelihood inequality proves the lemma; Appendix~\ref{app:likelihood} gives the time-uniform argument. When the split is at a training prefix the random variable $a(X)$ appears inside the sum. When the prefix is empty it takes the same deterministic value for every repetition of a frozen design and a fixed candidate. Uniformity over candidates is what later permits substituting the current fit into the errors of old experiments.

\paragraph{Local calibration.}
At a training prefix let $c_h(x)$ and $d_h(x)$ be the squared Hellinger distances between the two models for the following two test records. The current-core record contains the fresh uniform string and its observations. The one-step-plus-core record additionally includes the current observation and uniform action preceding the string. Each record is a projection of the complete logged suffix. Marginalization and the logged fair coin give $c_h+d_h\le2a_{\wh\theta,i}(x)$. With $\Delta q=\wh q-q^\st$ and $\Delta z=\wh z-z^\st$, the test-mass inequalities of Appendix~\ref{app:mass} give
\begin{align}
 \fused{\Delta q}^2&\le 8U c_h,\nonumber\\[-1pt]
 \Big(\sum_o\max_a\fused{\Delta z_{oa}}\Big)^2&\le 8AU d_h.
 \label{eq:drift}
\end{align}
To apply these bounds, insert $B_h^\st q_h^\st$ between $B_h^\st\wh q_h$ and $\wh B_h\wh q_h$. The first difference propagates $\Delta q$ through a full true continuation. The other is a difference of one-step joint predictions. By stability,
\[
 F_h(\wh q_h)
 \le \tfrac{\Lambda_B}{2}\Big(\fused{\Delta q}+
             \sum_o\max_a\fused{\Delta z_{oa}}\Big).
\]
The same right-hand side bounds $|g_v(\wh\phi_h-\phi_h^\st)|$ for every $v\in\K_h$. Consequently each squared quantity is at most $8Y a_{\wh\theta,i}(x)$. The initial-core experiment similarly gives $F_0^2\le2Y a_{\wh\theta,i}(\varnothing)$ and $F_0\le W$.

\subsection{The program given to the planner}
There are now three quantities at a prefix. The local operator error $F_h$ is defined using the truth. Its upper bound in terms of test-law discrepancies is also unknown, but Lemma~\ref{lem:mle} bounds the \emph{sum} of those discrepancies on past training data. The lifted rows and the set $\K_h$ are evaluable in the fit. The algorithm uses the known budget on the former quantities to constrain only the latter ones.
Set $S=8Y\beta$. At each refit, evaluate the histories in $\I_h$ using the current $\wh\theta$ and solve
\begin{align}
 b_h(x)&=\max_{v\in\K_h}g_v(\wh\phi_h(x)) \nonumber\\[-2pt]
 &\quad\text{subject to }\
 \sum_{x_i\in\I_h}g_v(\wh\phi_h(x_i))^2\le S .
 \label{eq:program}
\end{align}
The buffers are multisets. The initial scalar support is $b_0=\max\{s:|s|\le W,\ N_0s^2\le S\}$. Define the exact reward $b=b_0+\sum_{h=1}^H b_h(\tau_{h-1})$. Each stage selects its coefficient with only the prefix available at that stage.

\begin{proposition}[Stopping certificate]\label{prop:cert}
The reward satisfies $0\le b\le LW$. On the event of Lemma~\ref{lem:mle}, it satisfies \eqref{eq:intro} for every policy.
\end{proposition}
\begin{proof}
Summing local calibration over the disjoint training episodes gives $N_0F_0^2+\sum_{h,i}F_h(\wh q_h(x_i))^2\le S$. At any new prefix, \eqref{eq:global} puts its supporting coefficient in the feasible set of \eqref{eq:program}. Hence $F_h(\wh q_h(x))\le b_h(x)$ and $F_0\le b_0$. Apply \eqref{eq:telescope}. The zero coefficient is feasible and each stage objective is at most $W$.
\end{proof}

\paragraph{What is retained across fits.}
Consider two histories with identical lifted rows in the current fit. Every direction in \eqref{eq:program} evaluates to the same number on both. If that row has appeared $n$ times in the buffer, any feasible coefficient has absolute value at most $\sqrt{S/n}$ there. A new history can therefore benefit from several samples at other histories with the same row. Conversely, the type-one coefficient of \sect{sec:example} vanishes on all type-zero rows, leaving its value at a type-one query unconstrained by their number. Both conclusions follow from the same historical constraint. After refitting, these comparisons must use the new model's features on \emph{all} stored histories, since \eqref{eq:global} holds for one pair of models at a time.

The optimizer at a high-bonus prefix may differ from the local operator-error witness. To decide what to do there, we use the optimizer itself to choose a continuation and a score.

\section{From certificates to experiments}\label{sec:diagnostic}
The scalar maximum in \eqref{eq:program} determines whether to continue. Its optimizer determines what to observe. We derive the sampling rule from the same variables that implement the dual norm in the cone program.

\subsection{Implementing the maximizer}
For a single core action string the observation outcomes are disjoint. A signed coefficient on those outcomes can therefore be measured simply by executing the string and reporting a lookup-table score. An adaptive continuation creates a second possibility. A policy can prescribe one string after one observation and a different string after another. To sample such a policy it suffices to know how much action probability reaches each node of its test tree. The key convex step is that a coefficient in the dual ball can be split between these two kinds of experiments.

\begin{lemma}[A dual coefficient is an experiment]\label{lem:score}
At any stage, $\dual c\le b$ has a polynomial-size conic representation. A feasible representation yields an executable controller $K_c$, a mass $M_c\le\sqrt{U_h}\,b$ and an observable score $|f_c|\le1$ such that, for every valid state $q$,
\begin{equation}
 c^\top q=M_c\,\E_q^{K_c}f_c .
 \label{eq:scored}
\end{equation}
\end{lemma}
\begin{proof}
Suppress the stage index. Split $c=s+t$ with cost $b_1+b_2\le b$. For the group part impose $|s_j|\le d_{\alpha(j)}$, $d\ge0$, $\|d\|_2\le b_1$. For the second part place a nonnegative realization weight $p$ on the test-prefix tree, with root mass $b_2$. At an observation branch copy the incoming weight to each possible observation. At an action decision split it over actions:
\[
 p(\ell o)=p(\ell),\qquad
 \textstyle\sum_a p(\ell oa)=p(\ell o).
\]
Include each action at a decision node, allowing exits from the tree. Require $|t_j|\le p(j)$ on maximal leaves, and $t_j=0$ elsewhere. At a positive-mass decision the ratio $p(\ell oa)/p(\ell o)$ implements its policy, so $p(j)=b_2\pi(j)$.

To see exactness, the group body's support, for unit cost, is $N_{12}$ by groupwise maximization and Cauchy--Schwarz. The signed flow body's support is $N_\Pi$: choose a maximizing policy and leaf signs. Splitting a unit cost between the two bodies forms their convex hull, whose support is the maximum of the two support functions. By \eqref{eq:fused} this is precisely the fused-dual unit ball. There are $O(AHD)$ flow variables per coefficient block. The remaining history constraint in \eqref{eq:program} is one Euclidean norm bound on a vector of linear evaluations.

For the group experiment take $M_s=\sum_\alpha d_\alpha$. Draw string $\alpha$ with probability $d_\alpha/M_s$, apply it and return $s_j/d_\alpha$ on listed outcome $j$, or zero on an unlisted outcome. Its expected score times $M_s$ is $s^\top q$. For the flow experiment run the realization policy. On reaching a maximal leaf $j$, return $t_j/p(j)$, and return zero on exit. The probability of leaf $j$ is $\pi(j)q_j$, so this experiment has mean $t^\top q/b_2$. Mix the experiments in proportions $M_s:b_2$. Their total mass satisfies $M_s+b_2\le\sqrt{U_h}\,b_1+b_2$. Omit zero-weight branches and use the constant zero score when the total mass vanishes.
\end{proof}

\paragraph{Two diagnostic continuations.}
At the current prefix, fix a feasible optimizer $(u,w)$ and its conic decomposition by a deterministic evaluation rule. Let $M_0=W/2$. To measure the first term in $g_v$, run $K_u$ with score $M_u f_u$. For the second term, observe $o$, then choose an action by a distribution $\alpha(a\mid o,x)\ge M_{w_{oa}}/M_0$. Such a distribution exists by \eqref{eq:geometry} and Lemma~\ref{lem:score}. Put spare probability on a fixed action. Run $K_{w_{oa}}$ and report
\[
 Z_w=\frac{M_{w_{oa}}}{\alpha(a\mid o,x)}\,f_{w_{oa}} .
\]
A zero-mass term has zero score. Each diagnostic score lies in $[-M_0,M_0]$. By \eqref{eq:update}, its conditional mean is
\[
 \E_\theta[Z_w\mid x]=\sum_{o,a} w_{oa}^{\top}z_h^\theta(o,a;x).
\] The same frozen rule measures the fitted and the true values of the linear functional.

Write $D_h^u,D_h^w$ for following $\pi$ up to stage $h$ and then using the indicated diagnostic. Define a logged mixture
\begin{equation}
 \rho=\tfrac13\pi+\tfrac{1}{3H}\sum_{h=1}^H(D_h^u+D_h^w).
 \label{eq:rho}
\end{equation}
Each diagnostic follows $\pi$ to its selected stage, so its complete log records both the history reached and the outcome of the selected test. The full-policy component also records the unmodified trajectory.

\subsection{Transferring the fitted bonus to the real system}
Fix the data and a policy $\pi$. Even with identical physical core coordinates, replacing the fit by the truth changes two mathematical objects. The sampling measure on prefixes changes. At any prefix reached by both models, the values of the lifted coordinates can also change. There is a third issue specific to a data-dependent support: its historical constraint was imposed with fitted feature rows. We first transport that constraint on the stored prefixes. Then we bound the discrepancy of the selected scores on the current policy's prefixes.

\paragraph{Historical feasibility in true coordinates.}
To compare support functions, use the same stored prefixes in the true representation. Let $\cE=2LS+16Y\beta$. Define $\mathcal V_{\st,k}$ by $|v_0|\le W$, $v_h\in\K_h$, and the single constraint
\[
 N_0v_0^2+\sum_{h,i}g_{v_h}(\phi_h^\st(x_i))^2\le\cE .
\]
Let $b_{\st,k}(\tau)$ be the support of this set at the concatenated true features: maximize $v_0+\sum_h g_{v_h}(\phi_h^\st(\tau_{h-1}))$. From \eqref{eq:drift} every tuple feasible in the separate fitted programs belongs to this true set. Indeed its total fitted energy is at most $LS$ and its total squared drift is at most $8Y\beta$ by Lemma~\ref{lem:mle}. Apply $(a+b)^2\le2a^2+2b^2$. This inclusion is uniform over tuples, so it holds when each stage optimizer is selected at the queried prefix.

\paragraph{On-policy and intervened observations.}
If the planner always picks mode $b$ at the right station, trajectory data from that choice have no power to detect the wrong success probability for mode $a$ in \sect{sec:example}. The transition block of $v$ can nevertheless put a large weight on $a$. Running the $w$ diagnostic samples just such a counterfactual action. Each diagnostic log also records the roll-in under $\pi$, so the law of the complete experiment reflects changes in the probability of reaching its starting history.

\begin{proposition}[Diagnostic transfer]\label{prop:transfer}
For the attained numerical reward $\widetilde b_k$, any frozen $\pi$, and its controller \eqref{eq:rho}, on the calibration event,
\begin{equation}
 \E_{\wh\theta_k,\pi}\widetilde b_k
 \le \E_{\st,\pi} b_{\st,k}
       +3HW\,\TV(P_{\wh\theta_k}^\rho,P_\st^\rho).
 \label{eq:transfer}
\end{equation}
\end{proposition}
\begin{proof}
Write $t_\nu=\TV(P_{\wh\theta_k}^\nu,P_\st^\nu)$ for the distance between complete logged experiments. Freeze the deterministic rule $x\mapsto v_h(x)$ that attains each numerical lower value, including its diagnostic controller and score. For $s\in\{u,w\}$ let $Z_h^s$ be the score returned by $D_h^s$; it lies in $[-W/2,W/2]$. The defining conditional means of these scores, followed by averaging over the roll-in, give
\[
 \E_{\wh\theta_k,\pi}\widetilde b_k
 \le v_0+\sum_{h=1}^H\sum_{s\in\{u,w\}}
      \E_{\wh\theta_k,D_h^s} Z_h^s.
\]
At a true trajectory, the tuple chosen at its visited prefixes is feasible in $\mathcal V_{\st,k}$. Its value is therefore at most $b_{\st,k}$ pointwise, and hence
\[
 v_0+\sum_{h=1}^H\sum_{s\in\{u,w\}}
      \E_{\st,D_h^s} Z_h^s
 \le \E_{\st,\pi}b_{\st,k}.
\]
Because the same frozen score is evaluated on each pair of logged experiment laws and its range has width $W$, the difference of its expectations is at most $Wt_{D_h^s}$. Thus the transfer cost is at most $WS_D$, where $S_D=\sum_h(t_{D_h^u}+t_{D_h^w})$. The mixture label in \eqref{eq:rho} is observed, so $3t_\rho=t_\pi+S_D/H$. Consequently $WS_D\le3HWt_\rho$, as claimed.
\end{proof}

Put $\tau_0=\eps/8$ and $\eta=\tau_0/(6HW)$. A nonstopping round has fitted bonus value above $\tau_0$ with the stated simulation accuracy. Proposition~\ref{prop:transfer} then implies
\begin{equation}
 \E_{\st,\pi^k}b_{\st,k}\ge\tau_0/2
 \quad\text{or}\quad
 \TV(P_{\wh\theta_k}^{\rho_k},P_\st^{\rho_k})\ge\eta.
 \label{eq:either}
\end{equation}
The algorithm collects both batches. The proof can therefore charge the round to whichever inequality holds.

\section{Exploration progress and complexity}\label{sec:progress}
The alternatives in \eqref{eq:either} are statements about one frozen round. To bound the number of rounds, we need a quantity that persists after the model is replaced. There is no reason for a direction with large fitted support at one round to retain a comparable length in the next fitted model. The transfer inequality lets us do the counting entirely in the space of the true environment.

\subsection{One coordinate system for the full run}
At stage $h$ let $\mathcal Q_h$ consist of true predictive states reachable with positive probability. Its span has dimension $r_h\le r$. Fix any orthonormal coordinates for that span and choose $r_h$ elements of $\mathcal Q_h$ whose column matrix $X_h$ has largest absolute determinant in those coordinates. A maximizer exists since the finite horizon and alphabets give finitely many physical histories. Every reachable state can now be expressed as
\[
 q_h^\star(x)=X_h\chi_h(x),\qquad \|\chi_h(x)\|_\infty\le1 .
\]
By multilinearity, replacing basis column $j$ by the query changes the determinant by the factor $\chi_{h,j}(x)$. Maximality forces the coordinate bound. Choosing basis vectors from the reachable set has a second consequence. Any coefficient measuring a valid state with magnitude at most $W$ has coordinate value at most $W$ on each basis column.

Let $T_h^\star$ denote the linear map $q\mapsto(q,(B_h^\star(o,a)q)_{o,a})$. Then $\phi_h^\star(x)=T_h^\star X_h\chi_h(x)$. The lift may have many coordinates, yet its first block is $q$ and its reachable span has exactly the dimension of the original reachable span. For a stage coefficient $v$, its coordinate row is $(T_h^\star X_h)^\top v$. Its entries are the evaluations of $v$ on the lifted basis states. These facts are used twice below, once for an observed sample matrix and once for an expected sample matrix. Include the scalar initial block as $h=0$ with $\chi_0=1$.

\subsection{Progress on the training histories}
Define the empirical matrix from the stage-$h$ training rows collected before round $k$,
\[
 \Sigma_{k,h}=\lambda_a I+\sum_{t<k,i}\chi_h(x_{t,h,i})\chi_h(x_{t,h,i})^\top,
 \qquad \lambda_a=\frac{\cE}{rLW^2}.
\]
At the root use one row per initial-core episode. Let $\ell_{k,h}=\chi_h^\top\Sigma_{k,h}^{-1}\chi_h$ at the current prefix. For a feasible tuple in $\mathcal V_{\st,k}$, write $w_h$ for its coordinate coefficients. Its total historical energy is at most $\cE$ and its total ridge energy is at most $\lambda_a rLW^2=\cE$. Cauchy--Schwarz in the block-diagonal matrix of the $\Sigma_{k,h}$ yields
\begin{align}
 b_{\st,k}^2 &\le 2\cE \sum_{h=0}^H\ell_{k,h},\nonumber\\
 b_{\st,k}^2 &\le R_a^2\sum_{h=0}^H\min(1,\ell_{k,h}),
 \qquad R_a^2=\max\{2\cE,L^2W^2\}.
 \label{eq:clip}
\end{align}
For the second line use the first if all leverages are at most one, and the range $b_{\st,k}\le LW$ otherwise.

Under the first alternative in \eqref{eq:either}, Jensen's inequality makes the expected sum of clipped leverages at least $\tau_0^2/(4R_a^2)$. Each of the $m_k$ independent roll-ins at a given stage has exactly the required marginal. The batch size in Appendix~\ref{app:budget} makes the realized sum over the batch at least one with the allocated conditional probability. If $C_h$ is the new batch outer-product sum at stage $h$, then
\begin{equation}
 \sum_h\log\frac{\det(\Sigma_{k,h}+C_h)}{\det\Sigma_{k,h}}
 \ge\log\Big(1+\sum_h\tr(\Sigma_{k,h}^{-1}C_h)\Big)
 \ge\log 2 .
 \label{eq:detgrowth}
\end{equation}
The first inequality follows from $\det(I+M)\ge1+\tr M$ for $M\succeq0$ and multiplication across stages. The trace counts the unclipped leverage. Thus a batch can double the product of the determinants even when every individual row has leverage below one.

\subsection{Progress under the diagnostic policy}
For the second alternative, reverse the product expansion in \eqref{eq:telescope}. Now the prefix follows the true model and the remaining operators belong to the fit. Write $G_{k,h}$ and $G_{k,0}$ for \eqref{eq:F} with the two model roles exchanged. For any controller $\nu$,
\begin{equation}
 \TV(P_{\wh\theta_k}^\nu,P_\st^\nu)
 \le G_{k,0}+\sum_{h=1}^H \E_{\st,\nu}G_{k,h}(q_h^\st).
 \label{eq:reverse}
\end{equation}
A logged controller contributes the same causal action and label factors in both products. After summing the absolute differences, its marginalized action rule is a valid history policy. This establishes the same bound for the logged experiment.

\paragraph{Why intervene inside the diagnostic?}
Suppose a future fit agrees with an old diagnostic on its terminal score. That agreement alone controls one expected linear combination. The determinant argument needs a bound on a sum of \emph{squared local errors} for the future fit, evaluated at states reached by the old diagnostic. A single terminal mean need not control such a sum. The learner therefore treats the diagnostic itself as a roll-in policy and cuts it at a random stage. The uniform test at the cut measures the coordinates in the local decomposition regardless of what continuation the future fit may prefer. This is the reason for nesting a second experiment inside the first one.

Write $J_j(\nu)=\Hell^2(P_\st^{E_j(\nu)},P_{\wh\theta_k}^{E_j(\nu)})$. Condition on all data used to choose the fit and controller. The remaining probabilities refer to a fresh episode. A cut records its physical prefix, any controller labels already drawn, and its new random choices. It can therefore measure errors under an adaptive roll-in. The following comparison locates each local test inside that log.

\begin{lemma}[Which cut measures a test]\label{lem:cutprojection}
Let $c_h(x),d_h(x)$ be the current-core and one-step-plus-core conditional squared distances from \sect{sec:certificate}, and $e_H(x)$ the conditional squared Hellinger distance for the final observation. Expectations below use physical histories under the true system and $\nu$. For $h<H$, $\E d_h\le2J_h$. At the root $c_1\le J_0$. For $2\le h\le H$, $\E c_h\le2A J_{h-1}$. The same last bound holds for $\E e_H$ at $h=H\ge2$, and $e_1\le J_0$ when $H=1$.
\end{lemma}
\begin{proof}
For joint laws $P,Q$ of $(X,Z)$ put $a(x)=\Hell^2(P_{Z\mid x},Q_{Z\mid x})$. Cauchy--Schwarz gives
\[
 \big(1-\Hell^2(P,Q)\big)^2
 \le \E_{P_X}(1-a(X))^2 \le 1-\E_{P_X}a(X).
\]
Thus $\E_{P_X}a(X)\le2\Hell^2(P,Q)$, including for any proper conditional version on a $Q_X$-null prefix. Project the log of $E_h$ to the physical prefix $x$, the current observation, uniform action, fresh next-core string and its observed outcomes. Under each model the conditional law of this suffix is precisely the experiment defining $d_h$. The true prefix law is the roll-in law of $\nu$. Hellinger distance decreases under projection, proving the first claim.

To recover the current-core experiment for $h>1$, use the preceding cut. Write $x=yoa$ with $y=\tau_{h-2}$. The joint probability of a physical prefix and compatible past controller labels factors into the environment's observation probabilities and a controller factor common to both models. Summing over labels gives the marginal action kernel $\bar\nu(a\mid y,o)$, including when $\nu$ retains a private random choice across stages. The true physical prefix masses under the two designs satisfy
\begin{align*}
 d_\nu(yoa)&=d_\nu(y)p_\st(o\mid y)\bar\nu(a\mid y,o),\\
 d_{\rm cut}(yoa)&=d_\nu(y)p_\st(o\mid y)/A .
\end{align*}
In particular $d_\nu(x)\le A d_{\rm cut}(x)$ at every prefix. After the replaced action, cut $h-1$ draws its own core string independently. Project to $(x,\text{fresh string},\text{test observations})$. The conditional squared distance equals $c_h(x)$, giving
\[
 \sum_x d_\nu(x)c_h(x)
 \le A\sum_x d_{\rm cut}(x)c_h(x)\le2A J_{h-1}.
\]
At the last stage the same log contains $o_H$, so projecting to $(x,o_H)$ proves the bound for $e_H$. At the empty prefix the initial cut measures the core directly and, for $H=1$, records $o_1$ directly. These projections give the root bounds without a change of prefix law.
\end{proof}

\begin{lemma}[Error measured by a random cut]\label{lem:nested}
For each frozen controller $\nu$,
\[
 G_{k,0}^2+\sum_{h=1}^H\E_{\st,\nu}G_{k,h}(q_h^\st)^2
 \le 24Y\sum_{j=0}^{H-1}J_j(\nu).
\]
\end{lemma}
\begin{proof}
We apply the local B-error decomposition of \citet[Propositions D.1--D.2]{bstable} in the reverse direction. Insert $\wh B_h\wh q_h$ between $\wh B_h q_h^\st$ and $B_h^\st q_h^\st$. The first difference propagates $\Delta q$ through a fitted continuation. The second gives $\Delta z$. Stability, \eqref{eq:drift}, and $(s+t)^2\le2s^2+2t^2$ imply at $h<H$
\[
 G_{k,h}^2\le4\Lambda_B^2\big(Uc_h+AUd_h\big).
\]
At $h=H$, the joint block is a scalar observation probability, common across actions. Writing $t_H$ for the conditional TV of $o_H$ gives $G_{k,H}\le\Lambda_B\fused{\Delta q}/2+t_H$. Since $t_H^2\le2e_H$, this yields $G_{k,H}^2\le4\Lambda_B^2 U c_H+4e_H$. The root contributes at most $2\Lambda_B^2 U J_0$.

For $H\ge2$, Lemma~\ref{lem:cutprojection} bounds the stage-one expectation by $4YJ_0+8YJ_1$ and each later interior expectation by $8Y(J_{h-1}+J_h)$. The terminal expectation is at most $16YJ_{H-1}$. Summing charges costs at most $24Y$ per cut. For $H=1$, the root costs at most $2YJ_0$ and the terminal stage at most $8YJ_0$.
\end{proof}

Take $\Gamma=24YL$. Each previous nested batch supplies $n_N$ copies of the same empty-prefix risk for the current fit, with $j$ integrated uniformly in its logged law. Its squared Hellinger distance is $H^{-1}\sum_j J_j(\rho_t)$. By Lemmas~\ref{lem:mle} and \ref{lem:nested},
\begin{equation}
 \sum_{t<k}\Big[G_{k,0}^2+
  \sum_{h=1}^H\E_{\st,\rho_t}G_{k,h}(q_h^\st)^2\Big]
 \le\Gamma\beta/n_N\le\eta^2/128 .
 \label{eq:nestedenergy}
\end{equation}
In the $t$th summand of \eqref{eq:nestedenergy}, the controller $\rho_t$ is the one that actually collected the data, while $G_{k,h}$ measures the error of today's fit. The empty designated prefix integrates this error over the full roll-in law of $\rho_t$. This is the distribution whose covariance we retain across fits.

\subsection{From predictive disagreement to determinant growth}
For clarity, first consider a scalar true coordinate and a linear error $G(\chi)=|w\chi|$. If the current fit has squared error at most $b$ on the old distributions, then
\[
 w^2\sum_{t<k}\E_{\st,\rho_t}\chi^2\le b .
\]
A current distribution with large $\E_{\st,\rho_k}|w\chi|$ must have appreciable second moment compared with that old coverage. For a general PSR two details change: the error is a seminorm, whose supporting coefficient may vary with the query, and the error budget is shared across stages. The domination argument below handles both while keeping the covariance fixed.

Use $G_{k,h}(\chi)$ also for its pullback to the coordinate space, including the root seminorm $G_{k,0}|\chi_0|$. Its value at each basis vector $e_j$ is at most $W$: the fitted full continuation of the true basis state and the true one-step update followed by a fitted continuation each have norm at most $W$ before the factor $1/2$. With $\lambda_b=\eta^2/(128rLW^2)$, put
\begin{align*}
 \Xi_{k,h}&=\lambda_b I+\sum_{t<k}\E_{\st,\rho_t}\chi_h\chi_h^\top,
 \\
 a_{k,h}&=\lambda_b\sum_j G_{k,h}(e_j)^2+
       \sum_{t<k}\E_{\st,\rho_t}G_{k,h}(\chi_h)^2 .
\end{align*}
By \eqref{eq:nestedenergy}, $\sum_h a_{k,h}\le\eta^2/64$. For each current $\chi$, take a linear functional $w$ supporting the seminorm there. Global domination $|w^\top y|\le G_{k,h}(y)$ gives $w^\top\Xi_{k,h}w\le a_{k,h}$. In particular,
\[
 G_{k,h}(\chi)^2\le a_{k,h}\,\chi^\top\Xi_{k,h}^{-1}\chi.
\]
To compare old and current coverage, define the expected leverage of the current diagnostic,
\[
 d_{k,h}=\E_{\st,\rho_k}\big[\chi_h^\top\Xi_{k,h}^{-1}\chi_h\big].
\]
Jensen’s inequality gives $\E_{\st,\rho_k}G_{k,h}(\chi_h)\le\sqrt{a_{k,h}d_{k,h}}$. In the second case of \eqref{eq:either}, the reverse telescope then forces
\begin{align*}
 \eta
 &\le \sum_h\sqrt{a_{k,h}d_{k,h}}\\
 &\le\Big(\sum_h a_{k,h}\Big)^{1/2}
      \Big(\sum_h d_{k,h}\Big)^{1/2}
 \le \tfrac{\eta}{8}\Big(\sum_h d_{k,h}\Big)^{1/2}.
\end{align*}
It follows that $\sum_h d_{k,h}\ge64$. The trace corresponding to this number is the sum of $\tr(\Xi_{k,h}^{-1}\E_{\st,\rho_k}[\chi_h\chi_h^\top])$. Adding the new population covariances gives log-determinant growth at least $\log 2$ by \eqref{eq:detgrowth}. The empirical matrix counts individual observed training prefixes. One update of $\Xi$ adds \emph{one expected outer product per stage} for the current diagnostic. The $n_N$ repetitions have already been used to shrink the error budget in \eqref{eq:nestedenergy}; multiplying this covariance update by $n_N$ would count their effect twice. Both matrices serve only in the analysis.

\subsection{Termination and policy value}
Over $K$ completed nonstopping rounds both determinants are nondecreasing. Since $\|\chi_h\|_2^2\le r$, their total increases from their initial ridges imply
\begin{equation}
 K\log2\le rL\log(1+Km_Kr/\lambda_a)
              +rL\log(1+Kr/\lambda_b).
 \label{eq:cap}
\end{equation}
The right-hand side grows logarithmically in $K$, giving $\softO(rH)$ updates. Appendix~\ref{app:budget} turns this inequality into a deterministic cap and allocates failure probabilities. The batch cost can now be read off directly. To detect the second alternative we need resolution $\eta$ of order $\eps/(LW)$. Turning a squared-error allowance of order $\eta^2$ into the historical bound \eqref{eq:nestedenergy} requires
\[
 n_N=O(YL\beta/\eta^2)=O(YL^3W^2\beta/\eps^2).
\]
Here $YL$ comes from summing the adjacent-cut local-error bounds and uniform mixing over the cut; $L^2W^2$ comes from the diagnostic transfer. The training batch uses $\softO((YL^2\beta+L^3W^2)/\eps^2)$ episodes per round. Multiplying the larger scale by $\softO(rL)$ rounds gives the episode count in Theorem~\ref{thm:main}. Expanding $YW^2=\Lambda_B^4 AU^2$ identifies where the stability and action-string factors enter. Increasing the precision of each frozen batch affects this episode count while leaving the growth per charged round at $\log 2$.

At a stopping round, each stage's certified value is within $\zeta=\eps/(256L)$ of its exact support. With the planning and simulation tolerances of \sect{sec:algorithm}, we obtain
\[
 \sup_\pi\E_{\wh\theta_k,\pi}b_k
 \le 3\eps/16+2\eps/128+L\zeta=53\eps/256<\eps/4.
\]
Proposition~\ref{prop:cert} gives uniform trajectory TV at most $\eps/4$. Transferring the value of a true optimal policy to the fitted model and transferring the output policy back costs at most twice this quantity. The final planning error is at most $\eps/2$. This proves the policy claim in Theorem~\ref{thm:main}. The same stopping event supports any subsequent choice of observable reward in $[0,1]$, with one planning call for that reward.

\section{A specialization to observable states}\label{sec:specialcase}
The distinction between a predictive coordinate, a sampling design, and a reachable direction is explicit in a finite-horizon Markov system with $S_d$ observable states. Let the observation at each stage be the current state. For stages $1,\ldots,H$, use one singleton test for each possible state, all with the empty action string. Write $p^\theta_h(\cdot\mid o,a)$ for a transition distribution. Before the current observation, $q$ is the distribution of that observation given the past. A proper choice of operators for $h<H$ is
\[
 B_h^\theta(o,a)v = v_o\,p^\theta_h(\cdot\mid o,a),
 \qquad B_H^\theta(o,a)v=v_o .
\]
For a signed input, condition on the first state $o$. The remaining policy-weighted absolute suffix mass is $|v_o|$, because after the first state all trajectory probabilities are nonnegative and sum to one. Summing over $o$ yields $\pn{\B^\theta_{H:h}v}=\|v\|_1$. With one action-string group, $\fused v=\|v\|_1$. Thus the class satisfies $\Lambda_B=1$, $U=1$ and $D=S_d$.

The true rank at an interior stage is the dimension of the span of the transition distributions that can be reached there under some preceding state and action. If these distributions are mixtures of $d$ fixed distributions, this rank is at most $d$, even when each distribution assigns positive probability to all $S_d$ observations. For a finite realizable candidate class, the high-probability bounds are $\softO(rAH^4\beta/\eps^2)$ episodes and $\softO(rH)$ calls to each oracle. The history rows still have explicit state coordinates of size $S_d$. In the likelihood budget, a class that independently varies many transition probabilities can have a large $\log|\Theta|$.

For this choice of cores the dual ball is an infinity-norm ball. A coefficient vector $c$ is measured by observing the state and reporting $c_o$, or its normalized version if a unit score is required. A transition coefficient is measured by choosing the prescribed action and then scoring the next state. The general signed-test construction extends these elementary state observations to core tests of several lengths and to policies that branch on intermediate observations. In \sect{sec:example}, both types emit \emph{ready}. The longer stage-two core identifies the probability of type one by imposing mode $a$ and observing the final outcome.

\section{Discussion}
The episode budget and the oracle-call budget arise at different points in the argument. A training prefix is useful when it expands coverage in the true prediction space. A diagnostic episode is useful when it constrains a fitted model that disagrees with the truth under the current experiment. Both kinds of information accumulate across fits because their progress can be expressed in one fixed coordinate system. The sample cost of the second mechanism depends on the squared scale of its signed score. A sharper connection between diagnostic error and the variance of that score could reduce this cost while retaining the same stopping rule and update count.

\bibliographystyle{ACM-Reference-Format}
\bibliography{references}
\appendix
\section{Ordinary likelihood under adaptive tests}\label{app:likelihood}
We prove Lemma~\ref{lem:mle}. Write $p_i,q_i$ for the conditional joint masses of the recorded episode under truth and a fixed candidate, given the past $\mathcal F_{i-1}$. The controller and the split convention are fixed at this point. Let $a_i(x)=\Hell^2(p_i(\cdot\mid x),q_i(\cdot\mid x))$. On the support of the truth, integrating first over $Z_i$ gives
\begin{align*}
 &\E_\st\!\left[e^{a_i(X_i)}
 \sqrt{q_i(X_i,Z_i)/p_i(X_i,Z_i)}\mid\mathcal F_{i-1}\right]\\
 &\quad=\sum_x\sqrt{p_{i,X}(x)q_{i,X}(x)}
       \,e^{a_i(x)}(1-a_i(x))\le1 .
\end{align*}
Indeed $e^t(1-t)\le1$ on $[0,1]$ and the marginal affinity is at most one. The successive products form a nonnegative supermartingale of initial value one. Ville's inequality and a union bound over $\Theta$ bound its logarithm by $\log(3|\Theta|/\delta)$ for every candidate at every sample size, except on an event of probability $\delta/3$. At a refit, the sum of the log likelihood ratios of the selected model is at least $-\xi$. Substitution gives $\sum_{i\le n}a_i(X_i)\le\beta$. When $X_i$ is empty its summand is the squared Hellinger distance of the entire experiment induced by the frozen controller. This explains how complete-episode MLE supplies the two different error budgets used in the proof.

\section{The two product expansions}\label{app:product}
Fix a physical trajectory and abbreviate the realized matrices by $T_h=B_h^\st(o_h,a_h)$ and $S_h=\wh B_h(o_h,a_h)$. Write $T_{j:i}=T_j\cdots T_i$, with empty products equal to the identity, and put $q=q_1^\st$, $\hat q=\wh q_1$. The scalar difference in trajectory observation probabilities under the fixed actions has two expansions:
\begin{align*}
 T_{H:1}q-S_{H:1}\hat q
 &=T_{H:1}(q-\hat q)\\
 &\quad+\sum_{h=1}^H T_{H:h+1}(T_h-S_h)S_{h-1:1}\hat q\\
 &=S_{H:1}(q-\hat q)\\
 &\quad+\sum_{h=1}^H S_{H:h+1}(T_h-S_h)T_{h-1:1}q .
\end{align*}
Each sum is obtained by subtracting adjacent products whose dividing point moves one step. In the first expression, $S_{h-1:1}\hat q$ is the fitted unnormalized prefix state; in the second, $T_{h-1:1}q$ is its true counterpart. Multiplying by the action weights of a fixed policy and summing over physical trajectories turns the prefix factors into the expectations in \eqref{eq:telescope} and \eqref{eq:reverse}. At a given $(x,o,a)$ the absolute sum over the future is at most the policy norm of its suffix array. The factor $1/2$ in $F$ and $G$ converts the absolute sum into total variation.

For a logged controller, fix a physical trajectory $\tau$ and a compatible sequence of internal labels $\lambda$. Its joint mass factors into the environment's observation probability along $\tau$ and the controller's product $k_\nu(\tau,\lambda)$ of action and label probabilities. The kernel is identical in both models once $\tau,\lambda$ are fixed. Conditional on $\tau$, normalizing $k_\nu$ over compatible labels therefore gives a common label kernel. In particular, logging preserves the divergence of this complete physical experiment, while displaying the mixture label lets one write its divergence as the weighted sum of component divergences. At a cut the new kernel retains all realized choices and replaces future choices by its own test decisions; the same factorization applies.

\section{From test laws to feature errors}\label{app:mass}
Consider nonnegative test arrays $p,q$ that arise from core probabilities with total scale $m_p,m_q$. Within any group, each array sums to at most its scale. On maximal tests, the policy-weighted mass is likewise bounded by that scale. Factoring $p_j-q_j=(\sqrt{p_j}-\sqrt{q_j})(\sqrt{p_j}+\sqrt{q_j})$ and applying Cauchy--Schwarz to each group gives
\[
 \fused{p-q}^2
 \le 2(m_p+m_q)\sum_j(\sqrt{p_j}-\sqrt{q_j})^2 .
\]
For $N_{12}$ sum the squared group bounds. For $N_\Pi$ apply weighted Cauchy--Schwarz to each policy on the maximal leaves, and use $\pi(j)\le1$ in the squared-root sum. This also gives the bound when the core lists contain both a test and one of its extensions.

For a uniformly selected core string there are $U_h$ groups. The total squared-root sum of the listed entries is at most $2U_hc_h$: include the common string label in the experiment and retain the disjoint listed observation outcomes in each group. With scales one, this gives the first inequality of \eqref{eq:drift}. In the one-step experiment the scale for $z_{oa}$ is $p(o)$ or $q(o)$. For $s_{oa}$ denoting its sum of squared-root discrepancies,
\begin{align*}
 \Big(\sum_o\max_a\fused{\Delta z_{oa}}\Big)^2
 &\le 2\Big(\sum_o(p(o)+q(o))\Big)\sum_o\max_a s_{oa}\\
 &\le 4\sum_{o,a}s_{oa}\le8AU_{h+1}d_h .
\end{align*}
The last inequality uses the uniform action and subsequent uniform string label. This argument charges a single action factor. Applied to the initial uniform-core data it also bounds the squared error of the initial vector after a stable continuation.

\section{Batch sizes, stopping and oracle cost}\label{app:budget}
Recall $\tau_0=\eps/8$, $\eta=\tau_0/(12HW)$, $R_a^2=\max\{2\cE,L^2W^2\}$ and $\Gamma=24YL$. For a sufficiently large numerical constant $c$ use
\begin{equation}
 m_k=\left\lceil \frac{cR_a^2}{\tau_0^2}
        \log\frac{cL(k+1)^2}{\delta}\right\rceil,\qquad
 n_N=\left\lceil128\Gamma\beta/\eta^2\right\rceil .
 \label{eq:batches}
\end{equation}
Condition on the past at a frozen round. The independent clipped leverages in the training batch are in $[0,1]$. Under the first alternative their total mean is at least $m_k\tau_0^2/(4R_a^2)$. A multiplicative Chernoff bound makes their realized sum at least one except with probability $\delta/[c'(k+1)^2]$ for a suitable $c'$. This is a conditional bound for each past that satisfies the alternative. It can therefore be intersected with the simultaneous calibration event afterwards. Hoeffding's inequality gives the simulation accuracy using $O(L^2W^2\eps^{-2}\log(cL(k+1)^2/\delta))$ fitted-model episodes. Choose the constants to spend a total of at most $2\delta/3$ on simulation and batch tails.

The initial matrices are multiples of the identity. At each stage the accumulated empirical trace above the ridge is at most $Km_Kr$ and the added population trace at most $Kr$. Monotonicity, the eigenvalue bound by trace, and $r_h\le r$ give the right side of \eqref{eq:cap}. Compute $K_{\max}$ by doubling an integer until \eqref{eq:cap} first fails at that integer. Since $m_K$ grows as $\log(K+1)$, this takes $K_{\max}=\softO(rL)$ as the cap. For example the right side divided by $rL$ is bounded by a constant times $\log(C(K+1)^3)$ for $C$ polynomial in the scalar parameters and their reciprocals where applicable. Doubling past a large constant times $rL\log(C(rL+1)^3)$ then suffices. Return failure after $K_{\max}$ completed nonstopping rounds. On the joint event this contradicts \eqref{eq:cap}, and a stopping round must have occurred earlier.

Every nonstopping round uses $Lm_k$ training episodes and $n_N$ nested episodes. Substitution gives
\begin{align*}
 Lm_k&=\softO\big((\beta YL^2+L^3W^2)/\eps^2\big),\\
 n_N&=O(\beta YL^3W^2/\eps^2).
\end{align*}
The second term controls the sum up to logarithms. Multiplying by the cap and using $YW^2=\Lambda_B^4 AU^2$ proves the episode bound. Each fit is followed by one bonus-planning call, with one extra call for final reward planning. The extra $(LW)^2$ in the nested budget comes from the TV resolution $\eta$ in \eqref{eq:transfer}.

\section{Numerical evaluation of the support}\label{app:numeric}
A stage query returns a feasible coefficient, the conic decomposition from Lemma~\ref{lem:score}, and a certified lower value within $\zeta=\eps/(256L)$ of the exact support. Use a deterministic solver and tie convention throughout a round. If a certified lower objective is negative, use the zero solution. This construction gives $0\le\widetilde b_k\le b_k\le\widetilde b_k+L\zeta$; the scalar root has the same tolerance. A planner accepting rewards in $[0,1]$ receives $\widetilde b_k/(LW)$ with the correspondingly rescaled tolerance.

For precision bounds let $d$ be the lifted dimension and $n$ the number of training rows at the stage. Each coefficient vector has Euclidean norm at most $R_c=\Lambda_B\sqrt d$, since $\fused{e_j}=1$ on a coordinate vector. Every lifted row has norm at most $M=(A+1)U$ by the group mass bound. Since $\fused{v}\ge\|v\|_2$, the coefficient constraints contain a Euclidean ball of radius $\Lambda_B/(4\sqrt A)$. After the historical constraint, an inscribed radius is
\[
 \min\!\left\{\frac{\Lambda_B}{4\sqrt A},
   \frac{\sqrt S}{2\sqrt{\max(1,n)}\,M}\right\}.
\]
In the extended formulation choose zero signed coordinates, split each block's positive cost allowance between group and flow with slack, take small positive group radii, and split each action flow uniformly. In the affine space of the decomposition, root and flow equalities, this gives strictly feasible inequalities. The smallest positive leaf weight scales as $A^{-H}$ times an input-dependent block budget, with logarithmic precision cost $O(H\log A)$ plus input bit lengths. Standard weak optimization for second-order cones yields feasible interior approximations at polynomial bit cost; shrinking toward this interior point absorbs solver residuals. A fixed completion supplies the action probabilities on branches of zero weight.

For approximate model access let $e$ bound the Euclidean error of each reported feature row and query. Replace the radius $\sqrt S$ by $\sqrt S-\sqrt n R_c e$ on reported rows. The resulting coefficients are feasible on exact rows by the triangle inequality. Conversely an exactly feasible coefficient can be shrunk toward zero by a relative amount $O(\sqrt n R_c e/\sqrt S)$ to meet the inward constraint. The objective loss is bounded by this amount times $W$, plus $O(R_c e)$ for query evaluation. Taking
\[
 e\le c\,\zeta\min\!\left\{1/R_c,\,
 \frac{\sqrt S}{\sqrt{\max(1,n)}\,R_cW}\right\}
\]
for a sufficiently small universal $c$ leaves room for solver tolerance and certified downward rounding. The required bits grow logarithmically in $1/e$. Reencoding prefixes, evaluating rewards and simulating policies thus use polynomial additional computation.

\section{Discrete classes with upper brackets}\label{app:bracket}
Suppose a deterministic episode cap $T$ is available and each designated prefix has at most $C_{\rm pre}$ possible values, including recorded prefix labels. Set $\kappa=(4TC_{\rm pre})^{-1}$ and $s=\kappa/(4T)$. An upper bracket $b$ assigns each allowed controller $\nu$ a nonnegative array $b^\nu$. A bracket covers a candidate $\theta$ if for every such $\nu$, $b^\nu\ge p_\theta^\nu$ pointwise and $\sum_\omega(b^\nu-p_\theta^\nu)(\omega)\le s$. Let $M_s$ brackets cover the class, each covering at least one candidate.

\begin{corollary}
Theorem~\ref{thm:main} holds with $\beta=3\log(9M_s/\delta)+2\xi+2$ whenever $T$ dominates the deterministic budget produced by this choice.
\end{corollary}
\begin{proof}
It suffices to recover Lemma~\ref{lem:mle}. Write $\ell=\log(9M_s/\delta)$. For each bracket, normalize $b_i=b^{\nu_i}$ to $\bar q_i=b_i/\sum b_i$ and apply the conditional supermartingale to $\bar q_i$. For the bracket covering the fit, $\sum_{i\le n}\log(\bar q_i/p_i)\ge-\xi-Ts$ on the observed episodes, since $1\le\sum b_i\le1+s$. Thus its accumulated conditional squared Hellinger risk is at most $\ell+\xi/2+Ts/2$ uniformly over refits.

For the same bracket put $I_i=\mathbf1\{b_{i,X}(X_i)\le2\kappa\}$. The true conditional expectation of $(b_i/p_i)e^{I_i}$ is at most $c_b=1+s+2(e-1)\kappa C_{\rm pre}$, by summing bracket mass on the small prefixes. Apply Ville's inequality to the products divided by $c_b$ per step and union bound over brackets. Since $\sum\log(b_i/p_i)\ge-\xi$ for the fitted bracket, $\sum I_i\le\ell+\xi+T\log c_b\le\ell+\xi+1$. Every fitted prefix of probability below $\kappa$ is counted, since bracket marginal excess is at most $s$. At the remaining prefixes, writing the bracket as the fitted mass plus its nonnegative excess gives conditional TV between the two of at most $s/\kappa$. Squared Hellinger distance is bounded by TV. The squared triangle inequality bounds total conditional risk by twice the bracket risk, plus $2Ts/\kappa$ and the small-prefix count. This gives $3\ell+2\xi+2$. The two simultaneous events cost at most $2\delta/9$. Substitution in the finite-class argument proves the claim.
\end{proof}

\clearpage
\section{Choosing the coefficient geometry}
A class can have both a fused stability parameter $\Lambda_B$ and an $\ell_1$ continuation bound $\pn{\B^\theta_{H:h}v}\le\Lambda_1\|v\|_1$. The latter gives a second realization of the same local witnesses. At a prefix constrain $\|u\|_\infty\le\Lambda_1/2$ and $\sum_a\|w_{oa}\|_\infty\le\Lambda_1/2$ for every $o$. Each coefficient block $c$ can be measured by drawing a uniform action string and returning $U_h c_j$ on its listed outcome $j$, zero on exit. This score has range bounded in magnitude by $U_h\|c\|_\infty$. The transition controller uses action masses proportional to these bounds as in \sect{sec:diagnostic}. The stage score bound becomes $W_1=\Lambda_1U$.

The certificate budget $S$ follows from the same fused calibration of $F_h$, because stability in the two geometries applies to the same model operators. For a feasible coefficient in the new geometry, the squared historical feature drift can be larger. Using $\|v\|_1\le\sqrt U\fused v$ in \eqref{eq:drift} bounds this drift by $O(AW_1^2 a_{\wh\theta,i})$. The true energy allowance therefore becomes $O(\beta(LY+AW_1^2))$. The nested local-error budget remains $\Gamma=O(YL)$, while the transfer threshold scales as $\eps/(LW_1)$. The resulting training cost is absorbed by the nested batch since $Y\ge A$ and $W_1\ge1$ (one may take $\Lambda_1\ge1$). Choosing the geometry with the smaller known score range before data collection yields
\[
 T=\softO\!\left(\frac{r\Lambda_B^2 AU H^4\beta}{\eps^2}
       \min\{\Lambda_1^2 U^2,\Lambda_B^2 U\}\right).
\]
Each choice gives the same $\softO(rH)$ bound on calls to each oracle.

\section{A capped trajectory certificate}
The stagewise construction yields a bounded version of \eqref{eq:intro} along each policy. If $b_0,b_h$ dominate the errors of \eqref{eq:telescope} pointwise and $B=b_0+\sum_h b_h(\tau_{h-1})$, then
\[
 \TV(P_\st^\pi,P_{\wh\theta}^\pi)\le3\E_{\wh\theta,\pi}\min(B,1).
\]
For $b_0>1$ the claim follows directly. Otherwise let $T_c$ be the first stage at which the sum exceeds one, taking $H+1$ when it never does. On crossing trajectories the integral of the positive part of fitted minus true density is at most the fitted crossing probability. On the complement, expand the absolute density difference as in \eqref{eq:telescope}. At stage $h$ relax the restriction on the full path to the prefix event $h<T_c$, then sum out its future. The resulting integral is at most $2\E_{\wh\theta,\pi}[b_0+\sum_{h<T_c} b_h]$. Both the crossing indicator and this pre-crossing sum inside the expectation are at most $\min(B,1)$, giving the claim.

\end{document}
